% ============================================================================
% Learning from Big Data - Assignment 1
% LaTeX report template
% ============================================================================

\documentclass[11pt]{article}

% ---- Packages (from the Rmd header-includes, plus what a LaTeX report needs) ----
\usepackage[a4paper,margin=2.5cm]{geometry}
\usepackage{float}
\usepackage{booktabs}          % to thicken table lines
\usepackage{graphicx}
\usepackage{caption}
\usepackage{amsmath}
\usepackage{xcolor}
\usepackage{multirow}
\usepackage{hyperref}
\hypersetup{colorlinks=true, linkcolor=blue, citecolor=blue, urlcolor=blue}

% ---- Helper commands for the template scaffolding ----
% A prompt box reproducing a "QUESTION" from the Rmd
\newcommand{\question}[2]{%
  \par\medskip
  \noindent\textbf{QUESTION #1.} #2\par\medskip
}
% A placeholder the student must replace with their own answer
\newcommand{\answer}[1]{%
  \textcolor{blue}{\emph{[YOUR ANSWER HERE: #1]}}
}

\title{Learning from Big Data: Assignment 1\\ Prof. Gui Liberali}
\author{Group number: \dotfill}


\begin{document}
\maketitle

\section*{Introduction}
This file provides a LaTeX template for Assignment 1 of the Learning from Big Data
course. It was prepared to save you time
so that you can focus on the theory and technical parts of the methods seen in class.

You have received the dataset of reviews, the three dictionaries with the training set
of words for each topic, a list of stopwords, and a validation dataset containing
sentences classified by a panel of human judges. Run your analysis in R to produce your predictions, likelihoods, tables, and figures, then report and interpret them here. This template does not reproduce any R code -- it is
where you write up and discuss the results of the analysis you ran separately.

Sections marked \textbf{QUESTION} are the parts you are asked to answer. Text shown as
\answer{...} marks where you should insert your own writing, tables, or figures.

This report has the following structure:
\begin{enumerate}
  \item \textbf{General Guidelines}
  \item \textbf{Research Question}
  \item \textbf{Data aggregation and formatting}
  \item \textbf{Supervised Learning -- The Naive Bayes classifier (NBC)}
  \item \textbf{Supervised Learning -- Inspect the NBC performance}
  \item \textbf{Unsupervised Learning -- Predict using LDA}
  \item \textbf{Unsupervised Learning -- Predict using Word2Vec}
  \item \textbf{Analysis -- answering the research question}
  \item \textbf{OPTIONAL -- run and interpret the AFINN lexicon for sentiment}
  \item \textbf{Appendix}
\end{enumerate}

\textcolor{red}{Before going further:} make sure you have run your analysis in R
(Tutorial 0 provides installation instructions) and have the outputs (predictions,
likelihoods, tables, figures) ready to paste into this report.

% ============================================================================
\section{General Guidelines}
% ============================================================================
Page length. The template has 4 pages. You are allowed to add 7 to 10 pages, not
including the appendix. There is a limit of pages, but you have the possibility of
using appendices, which are not limited in number of pages. Use your pages wisely.
For example, having a table with 2 rows and 3 columns that uses 50\% (or even 25\%)
of a page is not an efficient use of space.

% ============================================================================
\section{Submission}
% ============================================================================
Make sure that the files you submit on canvas contain the following:
\begin{enumerate}
    \item The pdf of your report named \textbf{Assignment\_1\_GROUPNUMBER.pdf}
    \item Your code, that can be run as-is, which prints all tables/figures you use in your pdf report.
    \item If you use any data not on the github, upload it as well.
    \item The $Reviews\_Raw$ file containing the calculated posterior content and sentiment likelihoods. 
\end{enumerate}



% ============================================================================
\section{Research Question}
% ============================================================================
\question{I}{Present here the main research question you are going to answer with
your text analysis. You are free to choose the problem and change it until the last
minute before handing in your report. However, your question should not be so simple
that it does not require text analysis. For example, if your question can be answered
by reading two reviews, you do not need text analysis; all you need is 10 seconds to
read two reviews. Your question should not be so difficult that you cannot answer in
your report. Your question needs to be answered in these pages.}

\answer{}

% ============================================================================
\section{Data aggregation and formatting}
% ============================================================================
\question{II}{Decide on how to aggregate or structure your data. The data you
received is at the review level (i.e., each row is a review). However, the variables
in the data are very rich and allow you to use your creativity when designing your
research question. For example, there are timestamps, which allow you to aggregate
the data at the daily level or even hourly level. There is information on reviewers,
which allows you to inspect patterns of rating by reviewer. There is information on
studios, and more. Please explicitly indicate how you structured your dataset, and
what your motivation is for doing so. Even if you are using the data at the review
level, indicate how and why that is needed for your research question.}

\answer{}

% ============================================================================
\section{Supervised learning -- the Naive Bayes classifier}
% ============================================================================

\subsection{Likelihoods}
\question{III}{Create the content likelihoods based on the three lists of words
provided (storyline, acting, visual). Be explicit on the decisions you took in the
process, and why you made those decisions (e.g., which smoothing approach you used).}

\answer{}

\question{IV}{Locate a list of sentiment words that fits your research question. For
example, you may want to look just at positive and negative sentiment (hence two
dimensions), or you may want to look at other sentiment dimensions, such as specific
emotions (excitement, fear, etc.). TIP: Google will go a long way finding these, but
do check if there is a paper you can cite that uses your list.}

\answer{}


% ============================================================================
\section{Supervised Learning -- Inspect the NBC performance}
% ============================================================================
\subsection{Compute confusion matrix, precision and recall}
\question{V}{Compare the performance of your NBC implementation (for content)
against the judges' ground truth by building the confusion matrix and computing the
precision and accuracy scores. Do not forget to interpret your findings.}

\answer{}

\begin{table}[H]
  \centering
  \caption{Confusion matrix -- content classifier vs.\ judges' ground truth (example layout).}
  \label{tab:confusion-matrix}
  \footnotesize
  \begin{tabular}{llccc}
    \toprule
    & & \multicolumn{3}{c}{Predicted} \\
    \cmidrule(lr){3-5}
    & & Storyline & Acting & Sound/Visual \\
    \midrule
    \multirow{3}{*}{Actual} & Storyline     & \answer{} & \answer{} & \answer{} \\
                            & Acting        & \answer{} & \answer{} & \answer{} \\
                            & Sound/Visual  & \answer{} & \answer{} & \answer{} \\
    \bottomrule
  \end{tabular}
\end{table}

% ============================================================================
\section{Unsupervised Learning: Predict using LDA}
% ============================================================================
\question{VI}{Using LDA, predict a variable in your dataset (in the case of movie reviews, use movie box office). Tip: you can pass a list of
reviews to the LDA package in order to get the posterior probability that the
reviews are about each topic. If you pass them all in a single document, you will
not get review-specific vectors.}

\answer{}


% ============================================================================
\section{Unsupervised Learning: Predict using Word2Vec}
% ============================================================================
\question{VII}{Using Word2Vec, predict a variable in your dataset (in the case of movie reviews, use movie box office).\\
Tip 1: you can reduce the dimensionality of the Word2Vec output with PCA/Factor
analysis. This will save you computing time.\\
Tip 2: Word2Vec gives you word vectors. You can then compute the average of these
word vectors for all words in a review. This gives you a vector describing the
content of a review, which you can use as your constructed variable(s).}

\answer{}


% ============================================================================
\section{Analysis -- Use the constructed variables to answer your research question}
% ============================================================================
\question{VIII}{Now that you have constructed your NLP variables for sentiment and
content using supervised and unsupervised methods, use them to answer your original
research question.}

\answer{}

% ============================================================================
\section{OPTIONAL: run and interpret sentiment with the supervised learning AFINN lexicon}
% ============================================================================
\question{IX (optional)}{Using the AFINN code you received in the lecture, compute
the sentiment using the syuzhet package/lexicon. Compare the performance of your NBC
implementation (for sentiment) assuming that the AFINN classification is the ground
truth, then build the confusion matrix and compute precision and recall. Note that we
are now interested in understanding how much the two classifications differ and how,
but we are not implying that AFINN is error-free, far from it. We are interested in
uncovering sources of systematic differences that can be attributed to the algorithms
or lexicons. Do interpret your findings.}

\answer{}

\begin{table}[H]
  \centering
  \caption{Confusion matrix -- NBC sentiment vs.\ AFINN (example layout).}
  \label{tab:AFINN-confusion}
  \begin{tabular}{llcc}
    \toprule
    & & \multicolumn{2}{c}{Predicted (NBC)} \\
    \cmidrule(lr){3-4}
    & & Positive & Negative \\
    \midrule
    \multirow{2}{*}{AFINN} & Positive & \answer{} & \answer{} \\
                           & Negative & \answer{} & \answer{} \\
    \bottomrule
  \end{tabular}
\end{table}

% ============================================================================
\appendix
\section{Appendix}
% ============================================================================
Use this section (unlimited pages) for supplementary material: extra tables/figures,
robustness checks, additional prediction samples, or a link to your GitHub
repository in lieu of uploading your code separately.

\answer{}

\end{document}
